\section{Refining Discretization}
\label{sec:refinement}

Because the relaxation model relies on an aggregated time discretization $\discretization$, its optimal solutions may employ services that violate strict real-world constraints. In our algorithm sequence, the Refinement step is invoked whenever such an illegal lower-bound solution is produced. By systematically analyzing network flow, we identify specific chronologically invalid connections (which we denote as ``Impossible Departures''), services that too close to be physically realized (which denote as ``Inter Departures'') and super itineraries whose revenue contribution connot be achieved by any original itinerary combination (which we denote as ``Phantom Revenue''). We then strategically insert new time points to the discretization $\discretization$ to eliminate these infeasibilities. In this section, we first define the violation metrics for each of these three types of violations, and show that existence of any of these violations is necessary and sufficient for the relaxation solution to be illegal in the original model. Then, we present the refinement strategy to eliminate these violations, and prove that the algorithm will converge to a globally optimal solution after finite iterations. Note that the relaxation solutions can be illegal in the original model even when its objective value equals that produced by the upper bound, which means the algorithm is converged. We do not need to further refine the discretization to eliminate the violations in this case because we have already obtained a globally optimal solution.

\begin{definition}
\label{def:illegal_solution}
Let $\solution_{rel} = (\fleetAssignVar^*, \shareVar^*, \outsideShareVar^*)$ be an optimal solution to the relaxation model $\model{\discretization, \itinSetSuper}$ with an objective value of $Z_{rel}^*$. 

The relaxation solution $\solution_{rel}$ is defined as illegal with respect to the full model $\model{\discretizationFull}$ if there does not exist a passenger flow allocation $(\shareVar, \outsideShareVar)$ such that the combined assignment $\solution = (\fleetAssignVar^*, \shareVar, \outsideShareVar)$ is feasible in $\model{\discretizationFull}$ and achieves an objective value $Z = Z_{rel}^*$.
\end{definition}

\subsection{Impossible Departures}
The relaxation model may utilize resource units that physically have not arrived yet but are available in the model for the reason that we allow some short arcs in the relaxation model. This leads to infeasible service connections when we try to implement the solution in reality.
To detect unrealistic service connections, we apply a flow balance check at each active node $(i, t)$ in the time-space network.

At a specific node $(i, t)$, incoming flow from upstream must support outgoing flow. We categorize the incoming arcs of a resource type $f$ into components. Denote $\holdingArcIn(i, t, f)$ as the holding flow into node $(i, t)$ for resource type $f$, $\flightArc[-,legal](i, t, f)$ as services with feasible turn-around times, i.e., arcs $a = (i, t', j, t)$ where $t' + \travelTime_a \le t$, and $\flightArc[-,illegal](i, t, f)$ as short arcs where $t' + \travelTime_a > t$. We have
$$\inArcSet{i,t}=\holdingArcIn(i, t, f)\uplus \flightArc[-,legal](i, t, f) \uplus \flightArc[-,illegal](i, t, f)$$

The outgoing flows from node $(i, t)$ for resource type $f$ consist of service departures and holding flows. Denote $\departArc(i, t, f)$ as the departure flow from node $(i, t)$ for resource type $f$, and $\holdingArcOut(i, t, f)$ as the holding flow out of node $(i, t)$ for resource type $f$. We have
$$\outArcSet{i,t}=\departArc(i, t, f)\uplus \holdingArcOut(i, t, f)$$

At any node $(i, t)$, the available resource units for departure for resource type $f$ is the sum of holding flow and legal service arrivals, i.e., $\sum_{a\in\holdingArcIn(i, t, f)}\fleetAssignVar_{a,f} + \sum_{a\in\flightArc[-,legal](i, t, f)}\fleetAssignVar_{a,f}$. The actual departure demand is $\sum_{a\in\departArc(i, t, f)}\fleetAssignVar_{a,f}$. If the departure demand exceeds the available resource units, i.e.,
\begin{equation}
    \sum_{a\in\departArc(i, t, f)}\fleetAssignVar_{a,f} > \sum_{a\in\holdingArcIn(i, t, f)}\fleetAssignVar_{a,f} + \sum_{a\in\flightArc[-,legal](i, t, f)}\fleetAssignVar_{a,f}
\end{equation}
then we have impossible departures at node $(i, t)$ for resource type $f$. This means the model is relying on illegal service arcs to satisfy the departure demand, which is infeasible in reality.
For each active node $(i, t)$ and resource type $f$, we calculate the illegal departure quantity $\impDeparts(i,t,f)$:
\begin{equation}
    \impDeparts(i,t,f) = \sum_{a\in\departArc(i, t, f)}\fleetAssignVar_{a,f} - \sum_{a\in\holdingArcIn(i, t, f)}\fleetAssignVar_{a,f} - \sum_{a\in\flightArc[-,legal](i, t, f)}\fleetAssignVar_{a,f}
\end{equation}

Here, $\impDeparts(i,t,f)$ represents the net demand for resource units that cannot be satisfied by legal means. If $\impDeparts(i,t,f) > 0$, the model is forced to use illegal arrivals $\flightArc[-,illegal]$ to satisfy this departure demand. This metric directly quantifies the illegality.

By adding time points to the discretization $\discretization$ we can ensure that all these illegal departures will not occur in the refined model. Denote the refined discretization as $\discretization^{1}$. The refine procedure will be described in the Subsection~\ref{sec:refine_process}.
\subsection{Inter-Departure Violations}
In the relaxation model, multiple departures on the same segment may be scheduled within a segment's minimum inter-departure time, which violates operational separation constraints. To detect and quantify such violations we define the departure flow inside the window starting at time $t$ on segment $(i,j)$ and compare it with the allowable limit.

Let the set of arcs on segment $s = (i,j)$ be denoted by $\segArcSet{s}$. For each arc $a\in\segArcSet{s}$ write $a=(i,t_a,j,t'_a)$ so that $t_a$ denotes the departure time of $a$. Using the resource assignment variables $\fleetAssignVar_{a,f}$, we define the aggregated departure flow in the interval $[t, t+\minInterdepTime_{s})$ as
\begin{equation}
    \arcDepartFlow_{ijt} = \sum_{a\in\segArcSet{s}} \mathbf{1}\{t_a \in [t,\, t+\minInterdepTime_{s})\} \sum_{f\in\fleetSet} \fleetAssignVar_{a,f}.
\end{equation}

Here $\arcDepartFlow_{ijt}$ denotes the total number of departures measured in assigned flow on segment $(i,j)$ inside the minimum-separation window starting at $t$. In the full-model scope the instantaneous capacity permits at most one departure per separation window, hence the allowable upper bound is $1$. A separation violation occurs when the actual flow exceeds $1$.

We quantify the excess by the Separation Violation Flow
\begin{equation}
    \sepViolFlow_{ijt} = \max\bigl\{0,\; \arcDepartFlow_{ijt} - 1\bigr\}.
\end{equation}

Equivalently one may define the Separation Violation Amount
\begin{equation}
    \sepViolAmount_{ij}(t) = \arcDepartFlow_{ijt} - 1,
\end{equation}
and truncate to non-negative values when reporting violations.

Detection procedure: compute $\sepViolAmount_{ijt}$ for all segments $(i,j)$ and all discretized start times $t\in\discretization$. If there exists any $t$ such that
\begin{equation}
    \sepViolAmount_{ijt} > 0,
\end{equation}
then the solution contains an inter-departure violation (the second illegal condition in Theorem~\ref{thm:relax_illegal}). The metric identifies which segments and time windows require refinement or constraint tightening. By adding time points to the discretization $\discretization^{1}$ we can ensure that all these inter-departure violations will not occur in the refined model. Denote the refined discretization as $\discretization^{2}$. The refine procedure will be described in the Subsection~\ref{sec:refine_process}.
\subsection{Phantom Revenue}
When the relaxation solution $\solution = (\fleetAssignVar^*, \shareVar^*, \outsideShareVar^*)$ satisfies both resource availability and inter-departure constraints, it might still be illegal if it relies on "Phantom Revenue." This occurs when the revenue claimed by the aggregated super itineraries cannot be physically realized by any valid combination of their constituent original itineraries in the full discretization network.

Due to the highly coupled nature of the Sales-Based Linear Program (SBLP), the outside option share $\outsideShareVar$ acts as a dynamic variable. A local decrease in the attractiveness of available itineraries can shift demand to the outside option, thereby implicitly raising the choice bounds for other itineraries and reallocating passenger flows. Because of this endogeneity, evaluating Phantom Revenue via static slack analysis is insufficient. Instead, we formulate a Revenue Feasibility Subproblem (RFS) to establish a strict necessary and sufficient condition.

Given the fixed physical service assignment $\fleetAssignVar^*$, we seek to detect whether the target revenue contribution of each super itinerary $R$, defined as $Rev_{Rpl}^* = \demand_{mp} \fare_{Rl} \shareVar_{Rpl}^*$, can be attained by assigning valid shares $\shareVar_{rpl}$ to its constituent original itineraries $r \in \eqClassByPathSig_{(R)}$. We introduce continuous slack variables $s_{Rpl} \ge 0$ to capture the revenue shortfall for each super itinerary. The RFS is formulated as the following linear program:
\full{\begin{align}
RFS(\solution)=\min_{\shareVar, \outsideShareVar, s} \quad & \sum_{m\in \marketSet} \sum_{p\in \psgTypeSet} \sum_{l\in \fareClassSet} \sum_{R \in \itinSuperInMarket_m} s_{Rpl} \label{eq:rfs_obj}\\
\text{s.t.} \quad & \sum_{r \in \eqClassByPathSig_{(R)}} \demand_{mp} \fare_{rl} \shareVar_{rpl} + s_{Rpl} = \demand_{mp} \fare_{Rl} \shareVar_{Rpl}^* \quad \forall R \in \itinSuperInMarket_m, p\in \psgTypeSet, l\in \fareClassSet \label{eq:rfs_target}\\
& \outAttr_{mp} \shareVar_{rpl} \le \attr_{rpl} \outsideShareVar_{mp} \quad \forall m\in \marketSet, p\in \psgTypeSet, l\in \fareClassSet, r \in \itinInMarket_m \label{eq:rfs_ratio}\\
& \left(\frac{\adjOutAttr_{mp}}{\outAttr_{mp}}\right) \outsideShareVar_{mp} + \sum_{r \in \itinInMarket_m} \sum_{l\in \fareClassSet} \left(\frac{\adjAttr_{rpl}}{\attr_{rpl}}\right) \shareVar_{rpl} \le 1 \quad \forall m\in \marketSet, p\in \psgTypeSet \label{eq:rfs_limit}\\
& \sum_{m\in \marketSet}\sum_{p\in \psgTypeSet}\sum_{l\in \fareClassSet} \sum_{r \in \itinUsingArc{m}{a}} \demand_{mp} \shareVar_{rpl} \le \sum_{f\in \fleetSet} \fleetCap_f \fleetAssignVar_{af}^* \quad \forall a \in \flightArc \label{eq:rfs_cap}\\
& \shareVar_{rpl} \ge 0, \quad \outsideShareVar_{mp} \ge 0, \quad s_{Rpl} \ge 0 \nonumber
\end{align}}

Constraint (\ref{eq:rfs_target}) forces the underlying original itineraries to meet the relaxation revenue, using $s_{Rpl}$ to absorb any unavoidable deficit. Constraints (\ref{eq:rfs_ratio})-(\ref{eq:rfs_cap}) ensure that the recovered passenger flow strictly adheres to the physical capacity of $\fleetAssignVar^*$ and the original discrete choice bounds. 

The relaxation solution $\solution$ contains Phantom Revenue (and is therefore illegal) \textit{if and only if} the optimal objective value of the RFS is strictly greater than zero (i.e., $\sum s_{Rpl}^* > 0$). The time slots (i.e. a set of segment-time pairs) corresponding to the super itineraries with positive slack $s_{Rpl}^* > 0$ are the culprits that contribute to the Phantom Revenue. By adding time points to the discretization $\discretization^{2}$ we can ensure that all these phantom revenue violations will not occur in the refined model. Denote the refined discretization as $\discretization^{3}$. The refine procedure will be described in the Subsection~\ref{sec:refine_process}.
% ========================

The following theorem guarantees that the relaxation solution is illegal in the original model if and only if there are at least one of these three violations.

\begin{theorem}
\label{thm:relax_illegal}
A solution $\solution = (\fleetAssignVar^*, \shareVar^*, \outsideShareVar^*)$ to the relaxation model $\model{\discretization}$ is illegal in the original model $\model{\discretizationFull}$ if and only if one of the following three conditions holds:
\begin{enumerate}
    % \textbf{Impossible Departure}: 
    \item There exists at least one node $(i, t)$ and resource type $f$ such that the departure demand at $(i, t)$ for resource type $f$ exceeds the available supply, i.e., $\impDeparts(i,t,f) > 0$.
    % \textbf{Inter Departure}: 
    \item There exist a segment $(i,j)$ and departure time $t$ such that more than one service departs from node $i$ to $j$ within the minimum turn-around time, i.e., $\arcDepartFlow_{ijt} > 1$.
    % \textbf{Phantom Revenue}: 
    \item There exists a super itinerary in the solution whose revenue contribution cannot be achieved by any combination of original itineraries, i.e., $RFS(\solution) > 0$.
\end{enumerate}
\end{theorem}

% ========================

\subsection{Refinement Strategy}\label{sec:refine_process}
\paragraph{Impossible Departure}
If $\impDeparts(i, t, f) := \impDeparts > 0$, we have a violation. The culprits are the set of illegal arrival arcs $\mathcal{K} = \flightArc[-,illegal](i, t, f)$.

For each $k \in \mathcal{K}$, we determine its correction time $t'_{arr} = t_{\text{dep}}(k) + \travelTime_k$, which is the true physical arrival time.  We resolve the violation using the following greedy process. First, sort the culprits $\mathcal{K}$ by their flow volume $\fleetAssignVar_{k}$ in descending order. Then, greedily select culprits from $\mathcal{K}$ until the sum of their flow covers the deficit $\impDeparts$.
Finally, for each selected culprit $k$, add its true arrival time $t'_{arr}$ to the time discretization of node $i$ as $\discretization_i \gets \discretization_i \cup \{t'_{arr}(k)\}$. This ensures that in the next iteration, the model can no longer rely on the illegal arc $k$ to satisfy the departure demand at $(i, t)$, as it will now have a more accurate representation of the resource unit's availability time.

As a result, the time node grid is densified. In the next iteration, the previously problematic arc $k$ will be mapped to land at $t'_{arr}$ (which is later than $t$). This accurately delays the supply of the resource unit, rendering it unavailable for the impossible departure at $t$, thus the current relaxation solution will be infeasible in the next iteration.

\paragraph{Inter Departure}
If an inter-departure violation is detected on segment $(i,j)$ with $\sepViolAmount_{ij}(t) > 0$, it indicates that multiple physical services assigned to this segment are aggregated into an overlapping separation window. To resolve this, we identify the set of active original service arcs $k \in \segArcSet{s}$ whose physical departure times $t_a(k)$ fall within the violating interval $[t, t + \minInterdepTime_{s})$. 
Following a similar greedy logic to the impossible departures, we sort these culprit arcs by their flow volume $\fleetAssignVar_{k}$ and select the top arcs. For each selected arc $k$, we insert its true physical departure time $t_a(k)$ into the discretization of the origin node $i$ as $\discretization_i \gets \discretization_i \cup \{t_a(k)\}$. This densification ensures that the model can distinguish between these close departures in the next iteration, eventually forcing them to satisfy the minimum separation constraint.

\paragraph{Phantom Revenue}
When a super itinerary $R$ is identified as having phantom revenue (i.e., $s_{Rpl}^* > 0$ in the RFS), the current discretization is too coarse to represent the true physical connectivity or capacity constraints of its constituent itineraries $S_{(R)}$. To eliminate this infeasibility, we refine the discretization along the path of $R$. For each original itinerary $r \in S_{(R)}$ that contributes to the revenue shortfall, we extract the physical arrival and departure times of all its constituent service legs. Let $t^{phys}$ be such a physical time point at node $i$. We update the discretization $\discretization_i \gets \discretization_i \cup \{t^{phys}\}$. By splitting the time nodes at the transition points of the super itinerary, we break the equivalence class $S_{(R)}$, forcing the relaxation model to acknowledge the chronological or capacity conflicts that originally led to the revenue overestimation.

The following theorem guarantees that the refinement strategy will converge to a globally optimal solution after finite iterations.
\begin{theorem}
\label{thm:refine_converge}
第~\ref{sec:refine_process}小节中描述的细化策略，在松弛解于原始模型中非法时，将至少向时间离散化 $\discretization$ 中添加1个新时间点。因此，经过有限次迭代后，算法将收敛到全局最优解。
\end{theorem}